% Generated draft: requires mathematical and source review.
\documentclass[11pt]{article}
\usepackage[margin=25mm]{geometry}
\usepackage{fontspec}
\setmainfont{Latin Modern Roman}
\usepackage{amsmath,amssymb,mathtools,unicode-math}
\setmathfont{Latin Modern Math}
\usepackage{enumitem,needspace,xurl,hyperref}
\setlength{\parindent}{0pt}
\setlength{\parskip}{5pt}
\setlength{\emergencystretch}{3em}
\newcommand{\N}{\mathbb{N}}
\newcommand{\Z}{\mathbb{Z}}
\newcommand{\Q}{\mathbb{Q}}
\newcommand{\R}{\mathbb{R}}
\newcommand{\C}{\mathbb{C}}
\newcommand{\F}{\mathbb{F}}
\newcommand{\A}{\mathbb{A}}
\newcommand{\PP}{\mathbb{P}}
\newcommand{\E}{\mathbb{E}}
\newcommand{\eps}{\varepsilon}
\newcommand{\dd}{\,\mathrm{d}}
\begin{document}
\section*{UnsolvedMath ID: 30002022}
\textbf{Source problem code:} catalogue entry
\subsection{Definitions and mathematical statement}
The supplied record describes an integral algebra \(A\) whose complex module category is equivalent to the principal block \(\mathcal O_0\) of category \(\mathcal O\), and whose reduction over a field \(k\) models a corresponding modular category. These descriptions are supplied context: the integral coefficient ring, the construction of \(A\), the underlying Lie-theoretic data, the permitted modules, and the meaning of the modular model are not specified.

Write \(A_{\mathrm{aff}}\) provisionally for the affine analogue attributed in the record to Fiebig, and write \(\mathscr C_{\mathrm{aff}}\) for the category it controls. The clean statement calls this its module category, whereas the original statement only describes a category controlled by the affine algebra. Neither the algebra nor the precise category of modules is defined in the supplied material. The assertion that this analogue controls the full Lusztig conjecture is also supplied context; its hypotheses and precise mathematical meaning are absent.

Let \(\operatorname{Rep}(G_k)\) denote the rational representation category referred to in the record, and write \(\mathcal B_0(G_k)\) for its entire principal block, conventionally the block containing the trivial representation. The group \(G_k\), the characteristic and other hypotheses on \(k\), and the permitted dimensions of representations are not specified. The notation \(\operatorname{Rep}(G_k)\) is thus used for the ambient category, with \(\mathcal B_0(G_k)\) distinguishing the principal block intended by the question.

\textbf{Problem (permanent catalogue entry; source code catalogue entry).} Determine the relationship between \(\mathscr C_{\mathrm{aff}}\), the category controlled by the affine analogue of \(A\), and the entire principal block \(\mathcal B_0(G_k)\) of rational representations of \(G_k\).

This is a relationship question, not a supplied conjecture asserting an equivalence. The record does not specify a functor, an equivalence, a quotient description, or another precise comparison to establish. Consequently, a more specific mathematical assertion, with explicit domains and quantifiers, requires the missing constructions and hypotheses.
\subsection{Short English statement}
Problem 30002022 asks how the category associated with Fiebig's affine analogue of an integral algebra relates to the entire principal block of rational representations of \(G_k\). The supplied record does not define the affine algebra, specify the group and field hypotheses, or prescribe the type of categorical relationship sought. It therefore supports this qualified question, but not a precise equivalence conjecture.
\subsection{Sources}
\begin{itemize}
\item Supplied local UnsolvedMath registry record, permanent ID 30002022, title \emph{Affine Algebra Models for Principal Representation Blocks}, source problem code \texttt{catalogue entry}, in the Oberwolfach Reports Open Problems collection. The recorded proposed year is 2012; no individual proposer is supplied. The affine analogue is attributed to Fiebig in the supplied statement.
\item Supplied authoritative report citation: \emph{Representations of Finite Groups} (2012), Oberwolfach Reports, pp.\ 963--1019, DOI 10.4171/owr/2012/16. \url{https://doi.org/10.4171/owr/2012/16}. No problem-specific page or construction is reproduced in the supplied record.
\item Supplied literature source: \emph{Modular Koszul duality}. \url{https://arxiv.org/abs/1209.3760}.
\item Literature assessment dated exactly \texttt{2026-08-21}: the supplied curation labels the problem open and reports that the relationship to the whole principal block remains unresolved, while later work supplies structural dualities without the stated identification. This is supplied dated curation, not an independently verified assessment of the literature or current status. The background explicitly requests independent verification of time-sensitive status and cited claims.
\item The registry labels the statement \texttt{corrected\_verified} and reports comparison with the official Oberwolfach report. These are supplied verification claims; neither the report nor the literature source has been independently consulted for this draft.
\end{itemize}
\end{document}
